% =========================================================
\subsection{GeoAI Researcher / Model Analyst}
% =========================================================
For a GeoAI researcher, the questions are technical: \emph{how Geo-SMILE behaves across Hull},
and \emph{how faithful, spatially coherent and reproducible its explanations are}
(Figure~\ref{fig:stakeholder-researcher}).\par\medskip
% GeoAI Researcher / Model Analyst stakeholder card (Chapter 6). Same design as the other
% cards: two dashboard panels, four key figures, two text columns. City-wide: the card reports
% Hull-wide results of both branches and the explanation-quality evaluation, not one LSOA (the
% dashboard's Comprehensive mode places a selected LSOA in these distributions). Each figure
% appears once: the panels print rho_s and the map's colour cap, so tiles and text do not repeat
% them (checked by card_mode). Spatial clustering (LISA) and ATT Faithfulness are left to the
% dashboard's Comprehensive mode.
% Every image is the dashboard's own short-card researcher panel (Hull-wide, no area
% highlighted), drawn by dashboard.views.stakeholder_figures and exported at print size by
%   python export_dashboard_panels.py --stakeholder researcher_short
% then renamed from _E01033104_ to _hull_ (no LSOA is drawn).
%   01 signed coefficient against exact single-LSOA W1 shift, all 168 LSOAs (Geo branch)
%   02 |local coefficient| of the most important feature, six-month crime average (Feature branch).
% Text and tiles: generated by the dashboard LLM pipeline (generate_card_narrative.py
% --stakeholder researcher --openai-model gpt-4.1): master prompt geosmile-hull-v8 + researcher
% card-mode addendum. Technical wording and test statistics are allowed for this reader. The
% model marks figures and key terms itself. Prompts, evidence, raw output and checks:
% researcher_card_llm.json. The generated headline is not shown on the card.
%
% Preamble: \usepackage[most]{tcolorbox}, xcolor, graphicx.
% Copy into Chapter6/figs/: dashboard_researcher_short_hull_{1,2}.png and
% researcher_card_llm_content_openai.tex.
% =========================================================

\definecolor{cardInk}{HTML}{12171C}
\definecolor{cardMuted}{HTML}{57646D}
\definecolor{cardTeal}{HTML}{1F4E5F}
\definecolor{cardPink}{HTML}{FCE4EC}
\definecolor{cardWine}{HTML}{8A0033}
\definecolor{cardRule}{HTML}{D9DEE2}

% Same macros as planner_card.tex (\providecommand: defined once, whichever card comes first).
\providecommand{\hlpink}[1]{\begingroup\setlength{\fboxsep}{0.6pt}\colorbox{cardPink}{\textcolor{cardWine}{#1}}\endgroup}
\providecommand{\pv}[1]{\hlpink{#1}}
\providecommand{\kt}[1]{\hlpink{#1}}
\providecommand{\roletag}[1]{\tcbox[on line, colback=cardTeal, colframe=cardTeal, arc=2.5mm,
  boxsep=0pt, left=6pt, right=6pt, top=3pt, bottom=3pt]{\color{white}\sffamily\bfseries #1}}
\providecommand{\kicker}[1]{{\color{cardMuted}\sffamily\small\bfseries #1}}
\providecommand{\tile}[2]{\begin{tcolorbox}[enhanced, nobeforeafter, width=\cardtilewidth,
  equal height group=\cardtilegroup, colback=cardPink, frame hidden, arc=2mm,
  left=5pt, right=3pt, top=3pt, bottom=3pt, valign=top, before upper=\raggedright]
  {\bfseries\color{cardWine}#1}\par\vspace{1pt}
  {\sffamily\color{cardMuted}#2}\end{tcolorbox}}
\providecommand{\dashpanel}[3]{\kicker{#1}\par\vspace{2pt}%
  {\centering\includegraphics[width=#2\linewidth]{#3}\par}}

\def\cardtilegroup{researchertiles}
\def\cardtilewidth{0.2375\linewidth}
\providecommand{\cardHeadline}{[headline]}
\providecommand{\cardSecA}{[section 01]}
\providecommand{\cardSecB}{[section 02]}
\providecommand{\cardFooter}{[footer]}
\providecommand{\cardModel}{[model]}
\providecommand{\cardTileA}{\tile{--}{--}}
\providecommand{\cardTileB}{\tile{--}{--}}
\providecommand{\cardTileC}{\tile{--}{--}}
\providecommand{\cardTileD}{\tile{--}{--}}
\input{Chapter6/figs/researcher_card_llm_content_openai.tex}

\begin{figure}[!htbp]
\begin{tcolorbox}[enhanced, colback=white, colframe=cardRule, arc=3mm, boxrule=0.6pt,
                  left=3mm, right=3mm, top=2.5mm, bottom=2.5mm, before skip=0pt, after skip=2pt,
                  fontupper=\sffamily\normalsize\color{cardInk}]
\setlength{\parskip}{0pt}\raggedright%
\roletag{GEOAI RESEARCHER / MODEL ANALYST}\par\smallskip
{\color{cardMuted}How Geo-SMILE behaves across Hull, and how far its explanations can be trusted}\par\vspace{4pt}
% Sizing as on the police card: each panel larger, with its text section beside it.
% Row 1: Geo branch panel, with explanation behaviour beside it.
\begin{minipage}[t]{0.52\linewidth}\vspace{0pt}
  \dashpanel{01 \textperiodcentered{} Surrogate against exact masking}{1.0}{Chapter6/figs/dashboard_researcher_short_hull_1.png}
\end{minipage}\hfill
\begin{minipage}[t]{0.45\linewidth}\vspace{0pt}\raggedright
  \kicker{Explanation behaviour across Hull}\par\vspace{2pt}
  \cardSecA\par
\end{minipage}\par\vspace{6pt}
% Row 2: Feature branch panel, with fidelity, intervention and reproducibility beside it.
\begin{minipage}[t]{0.52\linewidth}\vspace{0pt}
  \dashpanel{02 \textperiodcentered{} Top feature, local effect}{0.9}{Chapter6/figs/dashboard_researcher_short_hull_2.png}
\end{minipage}\hfill
\begin{minipage}[t]{0.45\linewidth}\vspace{0pt}\raggedright
  \kicker{Fidelity, intervention and reproducibility}\par\vspace{2pt}
  \cardSecB\par
\end{minipage}\par\vspace{6pt}
% Row 3: explanation-quality figures.
\kicker{Explanation quality}\par\vspace{2pt}
\noindent\cardTileA\hfill\cardTileB\hfill\cardTileC\hfill\cardTileD\par\vspace{4pt}
{\color{cardRule}\rule{\linewidth}{0.4pt}}\par\vspace{1pt}
{\color{cardMuted}\footnotesize \cardFooter{} Panels are the dashboard's Hull-wide researcher view.
Text, tiles and \pv{highlights} generated by \cardModel{} (prompt geosmile-hull-v8, card
mode).\par}
\end{tcolorbox}
\caption{GeoAI Researcher / Model Analyst card: Hull-wide surrogate check (Geo branch), local
effect of the most important feature (Feature branch), explanation-quality figures and the
LLM-generated technical narrative.}
\label{fig:stakeholder-researcher}
\end{figure}
